\section*{Chapter \ref{ch:g2:measure} --- Money and Measures}

\begin{solution}{exo:g2:measure:1}
Pencil: cm. Classroom length: m. Feather: g. Schoolbag: kg. Water
bottle: L.
\end{solution}

\begin{solution}{exo:g2:measure:2}
$2$ m $= 200$ cm; \quad $300$ cm $= 3$ m; \quad $1$ kg $= 1000$ g.
\end{solution}

\begin{solution}{exo:g2:measure:3}
For example: $16 = 10 + 5 + 1$; \ $23 = 20 + 2 + 1$; \
$35 = 20 + 10 + 5$. (Other ways are possible.)
\end{solution}

\begin{solution}{exo:g2:measure:4}
$10 - 6 = 4$; \quad $20 - 17 = 3$; \quad $50 - 32 = 18$.
\end{solution}

\begin{solution}{exo:g2:measure:5}
Long hand on $6$, short between $4$ and $5$: half past four
($4{:}30$). Long hand on $3$, short just after $8$: quarter past
eight ($8{:}15$).
\end{solution}

\begin{solution}{exo:g2:measure:6}
Half past two: the long hand points down to the $6$, the short hand
sits between the $2$ and the $3$.
\end{solution}

\begin{solution}{exo:g2:measure:7}
After May: June. Before January: December. (Birthday month varies.)
\end{solution}

\begin{solution}{exo:g2:measure:8}
Spent: $3 + 4 = 7$. Change: $10 - 7 = 3$.
\end{solution}

\begin{solution}{exo:g2:measure:9}
$1$ kg $= 1000$ g, which is more than $800$ g: the sugar is
heavier, by $1000 - 800 = 200$ g.
\end{solution}

\begin{solution}{exo:g2:measure:10}
For example $20 + 10 + 5 + 5 = 40$ and $10 + 10 + 10 + 10 = 40$
(two ways with exactly four coins or bills).
\end{solution}
